\section{A planar word lemma with coefficients}
\label{planar:section}

This section proves the word-matching statement used in the rank argument.
The coefficient group may have torsion and need not be finitely presented.
The conclusion records equality of coefficient elements, as well as matching
letters. The proof follows the long-boundary principle of classical
small-cancellation theory associated with Greendlinger's lemma
\cite{greendlinger-1960}. We prove the coefficient-sensitive statement
in full.

Let $H$ be a group generated by a finite alphabet $S$, let $U$ be a disjoint
finite alphabet, and put $L=H*F(U)$.  Words will be written in
$S^{\pm1}\cup U^{\pm1}$.  An occurrence of a letter of $U^{\pm1}$ is called
an \emph{exterior occurrence}.  Between successive exterior occurrences
there is a word in $S^{\pm1}$; its value in $H$ is the corresponding
\emph{coefficient gap}.  For a cyclic word, this convention includes the
gap from the last exterior occurrence to the first.

A cyclic word is \emph{cyclically $H$-reduced} if no two successive
exterior occurrences are inverse letters separated by a coefficient gap
equal to $1$ in $H$.  An interval from one exterior occurrence through
another \emph{matches} a relator interval if their signed exterior letters
agree in order and all intervening coefficient gaps have equal values in
$H$.  The interval uses no position of the relator more than once.  A
\emph{full cyclic match} means that, after choosing cyclic starting
positions, the complete signed exterior lists have the same length and
agree, and all corresponding coefficient gaps, including the closing
gap, agree in $H$.

\begin{lemma}[Planar word lemma]
\label{planar:main}
Let $H$ be generated by a finite alphabet $S$, let $U$ be a disjoint
finite alphabet, and put $L=H*F(U)$. Let $\mathcal R$ be a family of
words in $L$ with the following properties.
Each $r\in\mathcal R$ has $m(r)\geq 7$ exterior occurrences, all positive
and with pairwise distinct underlying generators.  Distinct relators
share at most one generator from $U$.  Suppose that the natural map
\[
 H\longrightarrow D
 :=L/\langle\!\langle\mathcal R\rangle\!\rangle
\]
is injective.  If a cyclically $H$-reduced word $W$ has an exterior
occurrence and represents $1$ in $D$, then at least one of the following
conclusions holds:
\begin{enumerate}
\item $W$ has a full cyclic match with a relator or its inverse;
\item $W$ has two intervals disjoint in their exterior occurrences, each
  matching an interval of a cyclic relator or its inverse with at least
  $m(r)-3$ exterior occurrences, where $r$ is the relator matched by that
  interval.
\end{enumerate}
\end{lemma}

The two intervals in the second conclusion need not come from different
relators.  Disjointness refers to occurrences in $W$, and matching refers
to equality in $H$ of every gap strictly between the first and last
matched exterior letters.  In particular, no bound on the lengths of the
$S$-words spelling those gaps is asserted or needed.

\begin{proof}
We first turn a least-relator expression for $W$ into a planar graph.
The edges of this graph pair exterior occurrences.  Bigon faces give
consecutive matches, and a corner count will produce either one full
match or two long strings of bigons.

\paragraph{A finite filling over an arbitrary coefficient group.}
Choose an expression for $W$ in $L$ as a product of conjugates of words
from $\mathcal R^{\pm1}$ using the least possible number $N$ of those
words.  Free-product reduction shows that a cyclically $H$-reduced word
with an exterior occurrence is nontrivial in $L$, so $N>0$.

Give $H$ the presentation on $S$ with all words trivial in $H$ as
relations.  An equality in this presentation uses finitely many
relations: membership in their normal closure means a finite product
of conjugates.  Thus the chosen expression for $W$ can be witnessed
using a finite set of true $H$-relations.  Let $K_H$ be a finite
presentation complex for this finite set of relations, and wedge it
with one circle for each $u\in U$.  The resulting complex $K$ has a
natural homomorphism $\pi_1(K)\to L$.  We do not assume this homomorphism
is injective.

Realize the expression by a map from a sphere with $N+1$ open disks
removed to $K$.  Denote the distinguished boundary disk by $O$, and the
other disks by $B_1,\ldots,B_N$.  Read every disk boundary using its disk
orientation in the oriented sphere.  The boundary of $O$ reads $W^{-1}$,
and each $B_i$ reads one chosen relator or its inverse.  To construct the
map, cut the punctured sphere along disjoint spokes carrying the
conjugating paths.  Its remaining boundary reads the equality in $L$;
the finitely many chosen $H$-relations and free cancellations fill this
boundary.  Conversely, cutting any such mapped punctured sphere along
spokes expresses $W$ as a product of $N$ conjugates in $L$.

We may make the complex and map piecewise linear.  For example, realize
each relation with a mapping cylinder and a disk, triangulate, and use
the finite filling just described.  During later replacements we may
enlarge the finite set of true $H$-relations: any loop trivial in $L$
becomes fillable after adding the finitely many such relations used by
its equality.  Minimality throughout concerns $N$ in the actual group
$L$, not a possibly larger finitely presented approximation to it.

\paragraph{Tracks and connectedness.}
For every $u\in U$, choose a point $q_u$ in the interior of its circle
edge, avoiding the images of all vertices of a subdivision for the
piecewise linear map.  Near $q_u$ the target is an interval.  Its inverse
image consists of line segments across triangles, joined in pairs
across their sides.  These form disjoint cooriented arcs and circles in
the punctured sphere.  Each exterior occurrence on a boundary is the
endpoint of exactly one arc.  Each arc and each circle maps constantly
to its chosen point.  The two endpoints of an arc have opposite
crossing signs when read in the disk orientations.  This follows from
coorientation and the orientation of the surface; all disk-boundary
orientations are opposite the corresponding orientations as boundaries
of the punctured surface.

Regard the boundary disks as vertices and the arc tracks as edges of
an embedded graph $G$ on the sphere.  Closed tracks are not edges of
$G$; we retain them as part of the map.  We claim that $G$ is connected.
If it were not, choose a component $C$ not containing $O$ and a small
closed regular neighborhood of its disks and arc tracks.  This
neighborhood can be chosen disjoint from all other disks and tracks,
including the finitely many closed tracks.  Select the boundary curve
$\gamma$ facing the complementary component that contains $O$.
The disk side of $\gamma$ away from $O$ contains $C$ and at least one
inner disk; it may also contain other components, which causes no
problem.

The image of $\gamma$ avoids every $q_u$.  Cutting the exterior circles
at these points leaves a space retracting to $K_H$.  Consequently the
class of $\gamma$, after a change of basepoint, is an element of $H$
inside $L$.  Cap all relator holes on the side away from $O$ in a
presentation space for $D$.  The resulting mapped disk shows that this
element is trivial in $D$.  Injectivity of $H\to D$ makes it trivial in
$H$, hence $\gamma$ is trivial in $L$.  Replace that entire disk side
by a filling using only true $H$-relations and free cancellations.
At least one relator disk disappears, contradicting the choice of $N$.
Thus $G$ is connected.  This argument neither removes individual closed
tracks nor supposes that such tracks bound disks in the target.

\paragraph{Inner loops and dipoles.}
There is no graph loop at an inner vertex.  All exterior letters on
that disk have the same sign, whereas an arc pairs opposite signs.
Nor can an edge join two copies of the same relator.  Such copies would
have opposite orientations and meet the unique occurrence of the same
exterior generator $u$.  Base their boundary loops at the point $q_u$
on that occurrence.  Writing the positive occurrence as $u=ps$, a
positive based boundary has the form $sTp$, while the negative one is
$p^{-1}T^{-1}s^{-1}$; these are inverse loops.  The joining track is
constant and inserts no conjugating coefficient.  A small disk
neighborhood of the two disks and this track therefore has null
boundary in $L$.  Replacing its interior removes both relator disks,
again contradicting minimality.  Other tracks crossing the boundary of
this neighborhood impose no restriction on the replacement map.

For copies of different relators there is at most one joining edge:
the relators share at most one exterior generator, used once on each
disk.  Thus between any two inner vertices there is at most one edge.
Every inner vertex corresponding to $r$ has degree exactly $m(r)$.

\paragraph{Face equations and bigon matches.}
The complement of a small regular neighborhood of the connected
disk-and-arc graph consists of disks.  Their boundary curves follow the
face walks of $G$, even when a walk repeats vertices or edges.  Collapse
the exterior circles of $K$ to its basepoint and then take values in
$H$.  Each track becomes constant, and a sector between successive
track endpoints on a disk contributes its coefficient gap.  The map
on each complementary disk proves that the ordered product of its
sector values is $1$ in $H$.  Closed tracks inside that disk do not
affect this argument.

There are no monogon faces.  An inner loop has already been excluded;
a monogon at $O$ would pair two inverse exterior occurrences across a
gap trivial in $H$, contrary to cyclic $H$-reduction.  Any bigon incident
to an inner vertex consists of two edges joining that vertex to $O$.
Indeed, there are no parallel edges between inner vertices.  A face
walk of length two that traverses one edge twice would require degree
one at its inner endpoint, which is impossible since $m(r)\geq7$.

Such an inner--outer bigon matches consecutive exterior letters and
their intervening coefficient gap.  The exterior signs match because
$O$ is labelled $W^{-1}$ in its disk orientation.  Following $W$ along
the outer sector follows the corresponding inner relator in its disk
orientation.  The face equation, with this convention, is the inner
gap times the inverse of the outer gap equal to $1$.  It therefore gives
equality of the two elements of $H$, rather than equality merely up to
conjugacy.

We have now reduced the desired word matches to strings of bigons.
The remaining argument forces either a full cyclic string or two long
strings at different inner vertices.

\paragraph{The corner count.}
Give weight zero to every corner at $O$ and to every corner of a bigon.
For any other corner at an inner vertex, assign weight
\[
 \begin{cases}
  1/3,&\text{if both incident edges go to inner vertices},\\
  1/2,&\text{if exactly one incident edge goes to }O,\\
  1,&\text{if both incident edges go to }O.
 \end{cases}
\]
Occurrences are counted with multiplicity along a face walk.  A face
of length $\ell$ has total weight at most $\ell-2$.  For a bigon this
is equality.  For $\ell\geq3$ with no occurrence of $O$, the total is
$\ell/3\leq\ell-2$.  With exactly one occurrence of $O$, the two
neighboring corners have weight $1/2$ and the remaining $\ell-3$
corners weight $1/3$, again totaling $\ell/3$.  With at least two
occurrences of $O$, there are at most $\ell-2$ inner corners, each of
weight at most one.  These cases include repeated face walks and loops
at $O$.

Write $E$ and $F$ for the numbers of edges and faces of $G$, and $c(v)$
for the total corner weight at an inner vertex $v$.  Summing over faces
and using Euler's formula for the connected sphere graph with $N+1$
vertices gives
\begin{equation}
\label{planar:curvature}
 \sum_{v\ne O}\bigl(2-c(v)\bigr)
 \geq 2N-2E+2F=2.
\end{equation}

Call $2-c(v)$ the contribution of $v$.  A vertex of degree $m\geq7$
with no edge to $O$ has $c(v)=m/3>2$, so a positive vertex has an edge
to $O$.  At such a vertex, consider the cyclic gaps between successive
edges to $O$.  A gap that bounds a bigon costs zero.  A gap containing
no edge to an inner vertex, but not bounding a bigon, costs one.  A
gap containing $h\geq1$ edges to inner vertices costs
\[
 \frac12+\frac{h-1}{3}+\frac12
 =1+\frac{h-1}{3}.
\]
Each gap not bounding a bigon therefore costs at least one.

If every gap bounds a bigon, all edges at $v$ go to $O$.  Successive
edges at $v$ are also successive at $O$, in the opposite rotation.
They exhaust the cyclic rotation at $O$: otherwise a further half-edge
would interrupt one of these bigons.  Connectedness leaves no other
vertex.  The bigon equations then match every exterior occurrence of
$W$ with this relator, including all cyclic gaps.  This is the first
conclusion.

Otherwise positivity means that exactly one gap does not bound a
bigon, and that gap contains at most three edges to inner vertices.
Indeed, two such gaps cost at least two, while a gap containing
$h\geq4$ inner edges also costs at least two.  The other gaps form one
consecutive string of bigons joining at least $m-3$ edges to $O$.
Its bigon equations give a matching interval with at least $m-3$
exterior occurrences.  The contribution of this vertex is at most one.

Unless the first conclusion has already occurred, every positive
vertex contributes at most one.  The total contribution is at least
two, so there are at least two distinct positive vertices.  Take a
matching interval from each.  Each exterior occurrence of $W$ is the
endpoint of exactly one arc track, so intervals arising from different
inner vertices have disjoint exterior occurrences.  This proves the
second conclusion.
\end{proof}
